% =====================================================================
%  thmsmith Stage -1 proposal skeleton.
%  Written automatically by the thmsmith TS pipeline before Stage -1 runs.
%  Locked parameters: qid=eid\_proximal\_did\_bridge\_frontier, specialization=v1
%
%  HOW TO USE THIS FILE
%  --------------------
%  - The preamble (everything above the `% --- BEGIN CONTENT ---` marker)
%    and the `\section{...}` headers below are fixed by the pipeline.
%    Do NOT rewrite or rename them; the Step 3 reviewer subagent and the
%    downstream Stage 0 / Stage 1 prompts grep for these section titles.
% =====================================================================

\documentclass[11pt]{article}

\usepackage{amsmath,amssymb,amsthm,mathtools}
\usepackage{enumitem}
\usepackage[colorlinks=true,linkcolor=blue,citecolor=blue,urlcolor=blue]{hyperref}
\usepackage{natbib}

\newcommand{\E}{\mathbb{E}}
\renewcommand{\Pr}{\mathbb{P}}
\newcommand{\one}{\mathbf{1}}
\newcommand{\transpose}{\mathsf{T}}
\newcommand{\calH}{\mathcal{H}}
\newcommand{\calF}{\mathcal{F}}
\newcommand{\calR}{\mathcal{R}}
\newcommand{\R}{\mathbb{R}}
\newcommand{\plim}{\operatorname*{plim}}

\theoremstyle{definition}
\newtheorem{definition}{Definition}
\newtheorem{assumption}{Assumption}
\newtheorem{example}{Example}

\theoremstyle{plain}
\newtheorem{theorem}{Theorem}
\newtheorem{conjecture}{Conjecture}
\newtheorem{proposition}{Proposition}
\newtheorem{lemma}{Lemma}
\newtheorem{corollary}{Corollary}

\theoremstyle{remark}
\newtheorem{remark}{Remark}

\title{thmsmith Stage -1 proposal: \texttt{eid\_proximal\_did\_bridge\_frontier} / \texttt{v1}%
  \\ \large\textit{Finite negative-control bridge frontier for proximal difference-in-differences}}
\author{thmsmith Stage -1 proposer}
\date{\today}

\begin{document}
\maketitle

% --- BEGIN CONTENT ---  (everything above is fixed by the pipeline)

\paragraph{Locked parameters.}
This is an exact-identification proposal with locked tuple
\[
(\texttt{qid},\texttt{specialization},\text{SWIG},\text{treatment},
\text{target estimand},\text{identifying assumptions}).
\]
Here \(\texttt{qid}=\texttt{eid\_proximal\_did\_bridge\_frontier}\) and
\(\texttt{specialization}=\texttt{v1}\).  The SWIG is a two-group,
three-period panel with periods \(t\in\{-1,0,1\}\), group
\(G\in\{0,1\}\), treatment \(D_t=G\one\{t=1\}\), observed outcome
\(Y_t\), binary negative-control exposure \(W\in\{0,1\}\), binary
negative-control outcome \(Z\in\{0,1\}\), and latent trend state \(U\).
The treatment is the switch from \(D_0=0\) to \(D_1=1\) for the treated
group \(G=1\).  The target estimand is
\[
\theta_{\mathrm{ATT}}=\E\{Y_1(1)-Y_1(0)\mid G=1\}.
\]
The identifying assumptions are consistency; no anticipation through
period \(0\); negative-control exclusion for \((W,Z)\); preperiod bridge
stability linking the untreated trend bias from \(0\) to \(1\) to the
preperiod negative-control moment from \(-1\) to \(0\); positivity for
all displayed conditional probabilities; and, for the ordinary branch,
parallel trends.  The finite bridge-rank condition below is the candidate
frontier to be characterized, not an identifying assumption.  The focal
observable object is the finite pretrend-proxy bridge operator
\[
B_P(z,w)=\Pr_P(W=w\mid Z=z,G=1),\qquad
b_P(z)=\E_P(Y_0-Y_{-1}\mid Z=z,G=1)-\E_P(Y_0-Y_{-1}\mid Z=z,G=0),
\]
with bridge class
\[
\mathcal H_P=\{h\in\R^{\{0,1\}}:B_Ph=b_P\}.
\]
The proximal DiD adjustment is
\[
\Gamma_P(h)=\sum_{w\in\{0,1\}} h(w)\Pr_P(W=w\mid G=1),
\qquad
\theta_{\mathrm{prox}}(P,h)=\theta_{\mathrm{did}}(P)-\Gamma_P(h),
\]
where
\[
\theta_{\mathrm{did}}(P)=
\{\E_P(Y_1-Y_0\mid G=1)-\E_P(Y_1-Y_0\mid G=0)\}.
\]

\begin{abstract}
Difference-in-differences can be read as a negative-control design, but the published proximal and negative-control literatures do not isolate the finite preperiod object that decides when proxy information replaces parallel trends.  This proposal studies a two-group, three-period, binary-proxy design and asks for the exact frontier among three regimes: ordinary parallel trends identifies the ATT, a proximal bridge identifies the ATT after parallel trends fails, and neither restriction identifies the ATT.  The focal object is the observable matrix pair \((B_P,b_P)\) and the loading \(\pi_P(w)=\Pr_P(W=w\mid G=1)\), not a generic completeness condition.  The main conjecture states that, on a finite nondegenerate class, primitive negative-control exclusion and bridge stability imply proximal DiD identifies the ATT if and only if the bridge equation is solvable and \(\pi_P\) annihilates every null bridge direction, while ordinary DiD identifies it if and only if the untreated trend bias is zero.  A worked binary law computes all three regimes and shows that the parallel-trends and proximal-bridge restrictions are neither nested nor equivalent.  If resolved, the result would turn negative-control DiD diagnostics into an identification frontier rather than a placebo-test convention.
\end{abstract}

\section{Introduction}
Many applications have more credible negative controls than untreated-trend equality.  The question is whether two preperiods and a negative-control pair contain a finite bridge operator that identifies the treated counterfactual trend when ordinary parallel trends is false.  The headline claim is that the answer is governed by the observable pair \((B_P,b_P)\) and by a null-annihilator condition for the ATT loading.

\citet[\S3.2, pp. 1157--1160]{MiaoGengTchetgen2018ProxyVariables} identify causal effects through proxy rank conditions, and \citet[\S3, pp. 12--15]{SoferRichardsonColicinoSchwartzTchetgen2016NOC} interpret DiD as a negative-outcome-control strategy under constant unobserved-confounder association.  Direct negative-control DiD estimators in \citet{Gregori2025DIDNegativeControls} and calibration designs in \citet{ZhangEtAl2025NCCDiD} use negative controls to remove omitted-trend bias, but they do not state the finite pretrend-proxy rank frontier that separates ordinary DiD, proximal DiD, and nonidentification.  The kernel lives in exact identification, with panel notation used only to define the observed law.

The immediate consumer is negative-control DiD for real-world evidence and marketing panels, including \citet{PenningDeVriesGroenwold2023NegativeControls}, \citet{Gregori2025DIDNegativeControls}, and \citet{ZhangEtAl2025NCCDiD}.  A reader who accepts the conclusion would report the rank and null-loading frontier before treating negative-control balance as support for a DiD estimand.

\begin{itemize}[leftmargin=*]
\item Theorem 1: ordinary DiD and proximal DiD are two distinct exact-identification functionals in the finite binary-proxy design. Gap: \citet{SoferRichardsonColicinoSchwartzTchetgen2016NOC} gives the additive DiD control reading, while \citet{MiaoGengTchetgen2018ProxyVariables} gives generic proxy rank; neither states the joint frontier. Consumer: negative-control DiD applications in \citet{ZhangEtAl2025NCCDiD}.
\item Conjecture 1: the bridge frontier is necessary and sufficient for proximal ATT identification on a nondegenerate finite class. Gap: \citet[\S3.2, pp. 1157--1160]{MiaoGengTchetgen2018ProxyVariables} assumes a rank condition and \citet[\S2--\S3]{ParkRichardsonTchetgen2024SingleProxy} separates single-proxy ATT identification from two-proxy proximal identification; the conjecture gives the finite DiD null-loading frontier. Consumer: \citet{PenningDeVriesGroenwold2023NegativeControls}.
\item Theorem 2: a concrete \(2\times3\times2\times2\) witness exhibits parallel-trends-only, proximal-only, and neither-identified regimes, so the restrictions are neither nested nor equivalent. Gap: \citet{TchetgenParkRichardson2024UniversalDID} replaces additive parallel trends on outcome scale, and \citet{SoferRichardsonColicinoSchwartzTchetgen2016NOC} does not give pairwise non-nesting distributions. Consumer: empirical negative-control DiD design audits.
\end{itemize}

\section{Related work}
\citet[\S3.2, pp. 1157--1160]{MiaoGengTchetgen2018ProxyVariables} is the closest proxy-identification anchor; it uses two proxies and rank conditions for causal identification but does not specialize the rank object to finite preperiod DiD moments.  \citet[\S2--\S3]{TchetgenTchetgenYingCuiShiMiao2024Proximal} gives the broader proximal causal-inference framework in which outcome and treatment bridges replace unmeasured-confounding adjustment.  \citet[\S3--\S4]{CuiPuShiMiaoTchetgen2024JASA} develops proximal ATE and ATT efficiency theory under bridge existence and completeness; the present question precedes efficiency and asks when the finite DiD bridge functional is identified at all.  \citet[\S3, pp. 12--15]{SoferRichardsonColicinoSchwartzTchetgen2016NOC} is the closest ordinary-DiD negative-control anchor because it links DiD to equal unobserved-confounder associations over exposure groups and time.  \citet{Gregori2025DIDNegativeControls} is direct prior art: its doubly robust negative-control DiD estimator targets omitted-covariate bias using NCO/NCE pairs, while this draft asks when the finite population bridge object exists and is unique before estimation.  \citet{ZhangEtAl2025NCCDiD} is also direct prior art because it calibrates DiD with negative controls in real-world data; its calibration step is a consumer of, not a substitute for, the rank and nonidentification frontier in Conjecture~\ref{conj:frontier}.  \citet{PenningDeVriesGroenwold2023NegativeControls} compares COCA, DiD, generalized DiD, and double-negative-control methods and motivates a finite assumption-separation theorem.  \citet[\S2--\S3]{ParkRichardsonTchetgen2024SingleProxy} identifies ATT with a single negative-control outcome proxy under different restrictions, which makes the two-proxy bridge frontier non-redundant.  \citet{TchetgenParkRichardson2024UniversalDID} replaces additive parallel trends with odds-ratio equi-confounding for noncontinuous outcomes, which shows that scale and restriction choice matter.  \citet{RambachanRoth2023ParallelTrends} supplies the main DiD sensitivity comparator; it bounds violations using pretrend restrictions rather than negative-control bridge rank.  \citet{EgamiYamauchi2023MultiplePretreatmentDID} shows that multiple pretreatment periods can test and weaken trend assumptions, but its double-DID and K-DID moments are not proximal bridge equations.  \citet{ImbensKallusMao2021MinimalBridgePanel} treats panel fixed effects through bridge functions in factor settings and is adjacent to, but distinct from, the finite negative-control table studied here.  The in-catalogue prior art in \texttt{doc/API.md} records mature panel Q1 algebra and an empty ExactID catalogue; the closest banked proposal is \texttt{eid\_proximal\_phase\_v1}, whose reviewer asked for finite examples exhibiting all three proximal phases with explicit certificates.  This draft answers that banked gap in a DiD-specific observed law rather than by reusing the earlier generic bridge-envelope LP.

\section{Formal setup}
Let the observed law \(P\) be the distribution of
\[
O=(G,Y_{-1},Y_0,Y_1,W,Z),
\]
where \(G,W,Z\in\{0,1\}\), \(0<\Pr(G=1)<1\), and all conditional probabilities used below are positive.  The treatment path is \(D_t=G\one\{t=1\}\), so \(G=1\) indexes the treated group.  Potential outcomes are \(Y_t(d)\), with no anticipation \(Y_{-1}=Y_{-1}(0)\) and \(Y_0=Y_0(0)\).  The target is \(\theta_{\mathrm{ATT}}=\E\{Y_1(1)-Y_1(0)\mid G=1\}\).

Define the ordinary DiD estimand
\[
\theta_{\mathrm{did}}(P)=
\E(Y_1-Y_0\mid G=1)-\E(Y_1-Y_0\mid G=0).
\]
Define the untreated post-trend bias
\[
\Delta_P=\E\{Y_1(0)-Y_0(0)\mid G=1\}-\E\{Y_1(0)-Y_0(0)\mid G=0\}.
\]
Ordinary parallel trends is \(\Delta_P=0\).  The negative-control bridge uses
\[
B_P=
\begin{pmatrix}
\Pr(W=0\mid Z=0,G=1)&\Pr(W=1\mid Z=0,G=1)\\
\Pr(W=0\mid Z=1,G=1)&\Pr(W=1\mid Z=1,G=1)
\end{pmatrix},
\]
\[
b_P=
\begin{pmatrix}
\E(Y_0-Y_{-1}\mid Z=0,G=1)-\E(Y_0-Y_{-1}\mid Z=0,G=0)\\
\E(Y_0-Y_{-1}\mid Z=1,G=1)-\E(Y_0-Y_{-1}\mid Z=1,G=0)
\end{pmatrix},
\quad
\pi_P=
\begin{pmatrix}
\Pr(W=0\mid G=1)&\Pr(W=1\mid G=1)
\end{pmatrix}.
\]
The bridge fiber is \(\mathcal H_P=\{h\in\R^2:B_Ph=b_P\}\), and the proximal correction is \(\Gamma_P(h)=\pi_Ph\).  The finite bridge frontier is the triplet
\[
\mathfrak F(P)=\bigl(\mathbf 1\{\Delta_P=0\},\ \mathcal H_P,\ \pi_P\operatorname{ker}(B_P)\bigr).
\]

\[
\boxed{
\begin{gathered}
\texttt{qid}=\texttt{eid\_proximal\_did\_bridge\_frontier},\quad
\texttt{specialization}=\texttt{v1},\\
\text{SWIG}: (G,D_{-1},D_0,D_1,Y_{-1},Y_0,Y_1,W,Z,U),\quad
D_t=G\one\{t=1\},\\
\text{treatment}:D_1=1\text{ for }G=1,\quad
\text{target}:\theta_{\mathrm{ATT}}=\E\{Y_1(1)-Y_1(0)\mid G=1\},\\
\text{identifying assumptions}: \text{consistency, no anticipation, positivity, negative-control exclusion,}\\
\text{preperiod bridge stability; the candidate frontier is }(\mathcal H_P\neq\emptyset,\ \pi_P\operatorname{ker}(B_P)=\{0\}).
\end{gathered}}
\]

\section{Assumptions}
\begin{assumption}[Finite nondegenerate observed law]\label{ass:finite}
\(G,W,Z\in\{0,1\}\), \(Y_t\) is integrable for \(t=-1,0,1\), and every conditional probability entering \(B_P\) and \(\pi_P\) is strictly between zero and one.
\end{assumption}

\begin{assumption}[Consistency and no anticipation]\label{ass:consistency}
\(Y_t=Y_t(0)\) for \(t=-1,0\), and \(Y_1=G\,Y_1(1)+(1-G)Y_1(0)\).
\end{assumption}

\begin{assumption}[Ordinary parallel trends]\label{ass:parallel}
\(\Delta_P=0\).
\end{assumption}

\begin{assumption}[Negative-control exclusion]\label{ass:nc}
Conditional on the latent trend state \(U\) and group \(G\), \(W\) is not affected by the period-one treatment and \(Z\) has no direct causal effect on \(Y_t(d)\).  The pair \((W,Z)\) may be associated with \(U\).
\end{assumption}

\begin{assumption}[Preperiod bridge stability]\label{ass:stability}
If \(h\in\mathcal H_P\), then the untreated post-trend bias for the treated group equals \(\Gamma_P(h)\):
\[
\Delta_P=\pi_Ph.
\]
This is the finite DiD specialization of an outcome-bridge restriction; it is not assumed when ordinary parallel trends alone is invoked.
\end{assumption}

\begin{definition}[Candidate bridge frontier]\label{def:frontier}
The observable bridge frontier condition is
\[
\mathcal H_P\neq\emptyset
\quad\text{and}\quad
\pi_Pr=0\ \text{for every }r\in\operatorname{ker}(B_P).
\]
This condition is not assumed as identification.  Conjecture~\ref{conj:frontier} asserts that it is the primitive implication of negative-control exclusion, bridge stability, and finite support for the proximal branch.
\end{definition}

\begin{assumption}[Open-class realization]\label{ass:openclass}
The observed moments lie in a nonempty open subset of the finite observable simplex on which the displayed matrices and mean differences are generated by latent trend-state models satisfying Assumptions~\ref{ass:nc}--\ref{ass:stability}.  Boundary cases with zero denominators, deterministic \(W\), deterministic \(Z\), or collapsed group support are excluded.
\end{assumption}

\section{Main results}
\begin{theorem}[Theorem 1: two exact-identification functionals]\label{thm:functionals}
Under Assumptions~\ref{ass:finite}--\ref{ass:consistency}, ordinary parallel trends, Assumption~\ref{ass:parallel}, implies
\[
\theta_{\mathrm{ATT}}=\theta_{\mathrm{did}}(P).
\]
Under Assumptions~\ref{ass:finite}, \ref{ass:consistency}, \ref{ass:nc}, and \ref{ass:stability}, if Definition~\ref{def:frontier} holds, then every \(h\in\mathcal H_P\) gives the same value and
\[
\theta_{\mathrm{ATT}}=\theta_{\mathrm{did}}(P)-\pi_Ph.
\]
\end{theorem}
The first display is the usual DiD decomposition.  The second follows because bridge stability rewrites the untreated post-trend bias as \(\pi_Ph\), while Definition~\ref{def:frontier} makes \(\pi_Ph\) invariant over the bridge fiber.

\begin{conjecture}[Conjecture 1: finite bridge-rank ATT frontier (NOT YET PROVED)]\label{conj:frontier}
On the open finite class in Assumption~\ref{ass:openclass}, the following are necessary and sufficient for exact identification from \(P\).
\begin{enumerate}[label=(\roman*)]
\item Ordinary DiD identifies \(\theta_{\mathrm{ATT}}\) if and only if \(\Delta_P=0\).
\item Proximal negative-control DiD identifies \(\theta_{\mathrm{ATT}}\) through the finite pretrend bridge if and only if \(\mathcal H_P\neq\emptyset\) and \(\pi_Pr=0\) for every \(r\in\operatorname{ker}(B_P)\).
\item Neither ordinary DiD nor proximal DiD identifies \(\theta_{\mathrm{ATT}}\) if \(\Delta_P\neq0\) and either \(\mathcal H_P=\emptyset\) or there exists \(r\in\operatorname{ker}(B_P)\) with \(\pi_Pr\neq0\).
\end{enumerate}
Moreover each of the three regimes contains a nonempty open set of observable laws.
\end{conjecture}

\begin{theorem}[Theorem 2: binary-proxy non-equivalence witness]\label{thm:witness}
The three finite laws in Exhibit~9.1 satisfy Assumptions~\ref{ass:finite}--\ref{ass:consistency} and have no zero denominators.  The first is parallel-trends-only, the second is proximal-only, and the third is neither-identified.  Hence ordinary parallel trends and the finite proximal bridge frontier are neither nested nor equivalent as restrictions on \(P\).
\end{theorem}
The proof is direct matrix arithmetic.  The parallel-trends-only array has \(\Delta_P=0\) but \(\pi_P\operatorname{ker}(B_P)\neq\{0\}\), the proximal-only array has \(\Delta_P\neq0\) and invertible \(B_P\), and the third array has \(\Delta_P\neq0\) with \(b_P\notin\operatorname{range}(B_P)\).

\section{Sanity-check corollaries}
\begin{corollary}[No-proxy reduction]\label{cor:noproxy}
Under Assumptions~\ref{ass:finite}--\ref{ass:consistency} and \ref{ass:parallel}, Theorem~\ref{thm:functionals} collapses to standard two-period DiD because the ordinary branch sets \(\Delta_P=0\) and does not invoke any proxy correction.  If, in addition, the no-proxy specialization \(b_P=0\) and \(\pi_Pr=0\) for all \(r\in\operatorname{ker}(B_P)\) holds, then every admissible bridge has \(\Gamma_P=0\) and \(\theta_{\mathrm{prox}}=\theta_{\mathrm{did}}\).
\end{corollary}

\begin{corollary}[Full-rank bridge reduction]\label{cor:fullrank}
Under Assumptions~\ref{ass:finite}, \ref{ass:consistency}, \ref{ass:nc}, and \ref{ass:stability}, if \(\det(B_P)\neq0\), then the bridge solution is unique and the proximal ATT is
\[
\theta_{\mathrm{ATT}}=\theta_{\mathrm{did}}(P)-\pi_PB_P^{-1}b_P.
\]
This is the finite binary analogue of the usual proximal rank condition.
\end{corollary}

\paragraph{Exhibit 9.1: computed \(2\times3\) binary-proxy frontier witness.}
All three laws use \(\Pr(G=1)=1/2\), \(\Pr(Z=1\mid G=g)=1/2\), latent state \(U=Z\), and the displayed \(B_P(z,w)=\Pr(W=w\mid Z=z,G=1)\).  For \(G=0\), set \(\Pr(W=1\mid Z=z,G=0)=1/2\).  Conditional on \((G,U,W)\), the outcomes are degenerate at the displayed means, with \(Y_{-1}=0\), \(m_g(u)=\E(Y_0-Y_{-1}\mid U=u,G=g)\), \(q_g=\E\{Y_1(0)-Y_0(0)\mid G=g\}\), and treatment effect \(\tau=1\) for \(G=1\).  Thus \(Z\) is predictive of \(Y\) only through \(U\), the observed treated postperiod mean is \(\E(Y_1\mid G=1,Z=z,W=w)=m_1(z)+q_1+\tau\), and the latent completion records \(Y_1(0)=m_1(z)+q_1\).  The entries below compute every object in Section~6 from these finite primitives.
\[
\begin{array}{c|c|c|c|c|c}
\text{regime} & B_P & (m_0(0),m_0(1)) & (m_1(0),m_1(1)) & (q_0,q_1) & \pi_P\\ \hline
\text{parallel-trends-only} &
\begin{pmatrix}.5&.5\\.5&.5\end{pmatrix} &
(0,0) & (0,0) & (0,0) &
\begin{pmatrix}.8&.2\end{pmatrix}\\[1.1em]
\text{proximal-only} &
\begin{pmatrix}.8&.2\\.3&.7\end{pmatrix} &
(0,0) & (.10,-.05) & (0,.04) &
\begin{pmatrix}.6&.4\end{pmatrix}\\[1.1em]
\text{neither} &
\begin{pmatrix}.5&.5\\.5&.5\end{pmatrix} &
(0,0) & (.10,-.10) & (0,.06) &
\begin{pmatrix}.8&.2\end{pmatrix}
\end{array}
\]
These primitives imply \(b_P=m_1-m_0\), \(\Delta_P=q_1-q_0\), \(\theta_{\mathrm{ATT}}=1\), and \(\theta_{\mathrm{did}}=1+\Delta_P\).  For the first law, \(\operatorname{ker}(B_P)=\operatorname{span}\{(1,-1)^\transpose\}\) and \(\pi_P(1,-1)^\transpose=.6\), so bridge correction is not invariant even though \(\Delta_P=0\).  It is parallel-trends-only.  For the second law, \(\det(B_P)=.50\) and
\[
h=B_P^{-1}b_P=(.16,-.14)^\transpose,\qquad
\pi_Ph=.6(.16)+.4(-.14)=.04=\Delta_P,
\]
so \(\theta_{\mathrm{prox}}=(1+.04)-.04=1=\theta_{\mathrm{ATT}}\), while ordinary DiD is biased by \(.04\).  For the third law, \(b_P\notin\operatorname{range}(B_P)=\operatorname{span}\{(1,1)^\transpose\}\) and \(\Delta_P=.06\), so no bridge exists and ordinary DiD is biased.  The internal checks are immediate: each row of \(B_P\) sums to one, every entry lies in \((0,1)\), the determinant in the full-rank case is nonzero, the singular cases have an explicit null vector, and the outcomes above give a finite observed law plus a latent untreated completion without zero denominators or collapsed arms.

\section{Proof-sketch route}
For Theorem~\ref{thm:functionals}, the route is: ATT decomposition \(\leftarrow\) unfold Section~6 definitions \(\leftarrow\) substitute consistency and no anticipation \(\leftarrow\) isolate \(\Delta_P\) \(\leftarrow\) either set \(\Delta_P=0\) or use bridge stability \(\Delta_P=\pi_Ph\) \(\leftarrow\) use null-annihilation to show invariance over \(\mathcal H_P\).  For Conjecture~\ref{conj:frontier}, the route is: finite frontier headline \(\leftarrow\) unfold bridge fiber and ATT loading \(\leftarrow\) prove sufficiency by the preceding decomposition \(\leftarrow\) prove necessity by constructing two latent trend-state completions with the same \(P\) but distinct \(Y_1(0)\) means whenever \(b_P\) is infeasible or \(\pi_P\) moves along a null direction \(\leftarrow\) show the completions persist under small perturbations of \(P\) using finite-dimensional rank-stratum arguments \(\leftarrow\) verify the three explicit cells in Exhibit~9.1 and thicken them to open neighborhoods by determinant and range-continuity inequalities.  This chain has four substantive steps beyond definitions and assumptions: the ATT-bias decomposition, the bridge-null invariance lemma, the equal-\(P\) latent-completion necessity construction, and the open-neighborhood thickening.  For Theorem~\ref{thm:witness}, the route is the displayed matrix calculation plus the non-nesting logic across the three rows.

\section{Honest scope flags}
Theorem~\ref{thm:functionals}, Corollary~\ref{cor:noproxy}, Corollary~\ref{cor:fullrank}, and Theorem~\ref{thm:witness} are intended as finite algebraic statements.  Conjecture~\ref{conj:frontier} is not yet proved because the necessity direction requires a latent completion argument showing that bridge infeasibility or non-invariant null directions can change \(\theta_{\mathrm{ATT}}\) without changing \(P\).  The proposal does not claim efficiency theory, staggered-adoption inference, or validity for continuous proxies.  The open-class realization assumption may need to be narrowed if the latent completion requires additional support inequalities.

\section{Iteration trail}
\begin{itemize}[leftmargin=*]
\item v1 (cold-start) -- initial draft centered on the finite pretrend-proxy bridge operator, a binary witness, and non-equivalence of parallel trends and proximal bridge restrictions; prior verdict: none.
\item v2 (revise) -- removed the bridge frontier as an identifying assumption, added direct negative-control DiD comparators, replaced the witness with finite observed and latent laws, and fixed the no-proxy sanity reduction; prior verdict: REVISE.
\end{itemize}

\section*{Bibliography}
\begin{thebibliography}{99}
\bibitem[Cui et~al.(2024)]{CuiPuShiMiaoTchetgen2024JASA}
Cui, Y., H. Pu, X. Shi, W. Miao, and E. J. Tchetgen Tchetgen (2024).
Semiparametric proximal causal inference.
\textit{Journal of the American Statistical Association}.

\bibitem[Egami and Yamauchi(2023)]{EgamiYamauchi2023MultiplePretreatmentDID}
Egami, N. and S. Yamauchi (2023).
Using multiple pretreatment periods to improve difference-in-differences and related designs.
\textit{Political Analysis}.

\bibitem[Gregori(2025)]{Gregori2025DIDNegativeControls}
Gregori, P. (2025).
Difference-in-differences with negative controls.
\textit{Proceedings of the European Marketing Academy}.

\bibitem[Imbens et~al.(2021)]{ImbensKallusMao2021MinimalBridgePanel}
Imbens, G., N. Kallus, and X. Mao (2021).
Controlling for unobserved confounding in panel data using minimal bridge functions.
arXiv:2108.03849.

\bibitem[Miao et~al.(2018)]{MiaoGengTchetgen2018ProxyVariables}
Miao, W., Z. Geng, and E. J. Tchetgen Tchetgen (2018).
Identifying causal effects with proxy variables of an unmeasured confounder.
\textit{Biometrika} 105(4), 987--993.

\bibitem[Park et~al.(2024)]{ParkRichardsonTchetgen2024SingleProxy}
Park, C., D. B. Richardson, and E. J. Tchetgen Tchetgen (2024).
Single proxy control.
\textit{Biometrics}.

\bibitem[Penning de Vries and Groenwold(2023)]{PenningDeVriesGroenwold2023NegativeControls}
Penning de Vries, B. B. L. and R. H. H. Groenwold (2023).
A comparison of approaches to correct for unmeasured confounding using negative controls.
\textit{Statistical Methods in Medical Research}.

\bibitem[Rambachan and Roth(2023)]{RambachanRoth2023ParallelTrends}
Rambachan, A. and J. Roth (2023).
A more credible approach to parallel trends.
\textit{Review of Economic Studies}.

\bibitem[Sofer et~al.(2016)]{SoferRichardsonColicinoSchwartzTchetgen2016NOC}
Sofer, T., D. B. Richardson, E. Colicino, J. Schwartz, and E. J. Tchetgen Tchetgen (2016).
On negative outcome control of unobserved confounding as a generalization of difference-in-differences.
\textit{Statistical Science} 31(3), 348--361.

\bibitem[Tchetgen Tchetgen et~al.(2024a)]{TchetgenTchetgenYingCuiShiMiao2024Proximal}
Tchetgen Tchetgen, E. J., A. Ying, Y. Cui, X. Shi, and W. Miao (2024).
An introduction to proximal causal learning.
\textit{Statistical Science}.

\bibitem[Tchetgen Tchetgen et~al.(2024b)]{TchetgenParkRichardson2024UniversalDID}
Tchetgen Tchetgen, E. J., C. Park, and D. B. Richardson (2024).
Universal difference-in-differences.
\textit{Epidemiology}.

\bibitem[Zhang et~al.(2025)]{ZhangEtAl2025NCCDiD}
Zhang, Y. et~al. (2025).
Negative-control-calibrated difference-in-differences for real-world data.
\textit{npj Digital Medicine}.
\end{thebibliography}

\end{document}
